% answers.tex —— 1985 年·上海卷·理工农医类 · 参考答案与详细解析片段
% 由 tools/generate_projects.py 依来源总账自 1985/1985_数学真题_上海理_试题.md + 1985/1985_数学真题_上海理_解析正文.md + 1985/1985_数学真题_上海理_答案要点.md 生成，勿手改（改源 MD 后重跑）


\begin{gktip}
固定源为两页题面，不含答案。下列逐题推导均为本库独立复算；第 2 题第（5）小题保留固定源字面“$x=x^3$”与选项不一致的问题，不将选项倒推为题面。
\end{gktip}

\begin{answer}
\gksolhead{1}{}

\textbf{题目}：

（本题满分 30 分）

本题共有 10 个小题，每一个小题满分 3 分，只要求直接填写结果。

\Needspace{5\baselineskip}
（1） 不等式

\nopagebreak[4]
\[
|x+2|<1
\]

的解集是　　　　。

\Needspace{5\baselineskip}
（2） 函数

\nopagebreak[4]
\[
y=\sqrt[8]{x}-1
\]

的反函数是　　　　。

（3） 点 $(0,1)$ 到直线 $x+y=2$ 的距离是　　　　。

\Needspace{5\baselineskip}
（4） 函数

\nopagebreak[4]
\[
y=\sin\frac{\pi x}{2}
\]

的最小正周期是　　　　。

（5） 如果一个圆的圆心在点 $(2,4)$，并且经过点 $(0,3)$，那么这个圆的方程是　　　　。

（6） 若平面 $\alpha$ 及这个平面外的一条直线 $l$ 同时垂直于直线 $m$，则直线 $l$ 和平面 $\alpha$ 的位置关系是　　　　。

\Needspace{5\baselineskip}
（7） 函数

\nopagebreak[4]
\[
y=\sin^2x-\cos^2x
\]

的最小值是　　　　。

\Needspace{5\baselineskip}
（8） 方程

\nopagebreak[4]
\[
9^x-7\cdot3^x-18=0
\]

的解是　　　　。

（9） 若 $\operatorname{tg}\alpha=-2$，且 $\sin\alpha<0$，则 $\cos\alpha$ 的值是　　　　。

（10） 已知 $O$ 为直角坐标系的原点，点 $A$ 的坐标为 $(1,3)$，点 $B$ 的坐标为 $(0,3)$，若把 $\triangle OAB$ 绕 $y$ 轴旋转一周，则所得旋转体的体积是　　　　。

\Needspace{4\baselineskip}
\textbf{解析}

\gksubpart{（1）}

\Needspace{5\baselineskip}
由 $|x+2|<1$ 得 $-1<x+2<1$，所以解集为

\nopagebreak[4]
\[
(-3,-1).
\]

\gksubpart{（2）}

\Needspace{5\baselineskip}
由 $y=\sqrt[8]{x}-1$ 得 $x=(y+1)^8$。交换 $x\text{、}y$，并保留原函数值域条件，反函数为

\nopagebreak[4]
\[
y=(x+1)^8,\ x\ge-1.
\]

\gksubpart{（3）}

\Needspace{5\baselineskip}
点 $(0,1)$ 到直线 $x+y-2=0$ 的距离为

\nopagebreak[4]
\[
\frac{|0+1-2|}{\sqrt{1^2+1^2}}=\frac{\sqrt2}{2}.
\]

\gksubpart{（4）}

\Needspace{5\baselineskip}
$\sin u$ 的最小正周期为 $2\pi$。令 $u=\frac\pi2x$，得到

\nopagebreak[4]
\[
T=\frac{2\pi}{\pi/2}=4.
\]

\gksubpart{（5）}

\Needspace{5\baselineskip}
半径平方为

\nopagebreak[4]
\[
r^2=(2-0)^2+(4-3)^2=5,
\]

\Needspace{5\baselineskip}
所以圆方程为

\nopagebreak[4]
\[
(x-2)^2+(y-4)^2=5.
\]

\gksubpart{（6）}

\Needspace{5\baselineskip}
平面 $\alpha\perp m$，说明 $m$ 的方向是平面 $\alpha$ 的法向方向；直线 $l\perp m$，且题面明确 $l$ 在平面 $\alpha$ 外。因此 $l$ 与平面 $\alpha$ 平行，即

\nopagebreak[4]
\[
l\parallel\alpha.
\]

\gksubpart{（7）}

\nopagebreak[4]
\[
y=\sin^2x-\cos^2x=-\cos2x.
\]

因为 $-\cos2x$ 的最小值为 $-1$，故答案为 $-1$。

\gksubpart{（8）}

\Needspace{5\baselineskip}
令 $t=3^x>0$，则

\nopagebreak[4]
\[
t^2-7t-18=0,
\]

即 $(t-9)(t+2)=0$。由 $t>0$ 得 $t=9=3^2$，所以 $x=2$。

\gksubpart{（9）}

\Needspace{5\baselineskip}
由 $\tan\alpha=-2$ 且 $\sin\alpha<0$，可知 $\cos\alpha>0$。于是

\nopagebreak[4]
\[
\cos\alpha=\frac1{\sqrt{1+\tan^2\alpha}}=\frac1{\sqrt5}=\frac{\sqrt5}{5}.
\]

\gksubpart{（10）}

\Needspace{5\baselineskip}
三角形 $OAB$ 绕 $y$ 轴旋转形成底面半径为 $1$、高为 $3$ 的圆锥，因此

\nopagebreak[4]
\[
V=\frac13\pi\cdot1^2\cdot3=\pi.
\]

\gkanswer{答案（1）：$(-3,-1)$}

\gkanswer{答案（2）：$y=(x+1)^8\ (x\ge-1)$}

\gkanswer{答案（3）：$\dfrac{\sqrt2}{2}$}

\gkanswer{答案（4）：$4$}

\gkanswer{答案（5）：$(x-2)^2+(y-4)^2=5$}

\gkanswer{答案（6）：$l\parallel\alpha$}

\gkanswer{答案（7）：$-1$}

\gkanswer{答案（8）：$x=2$}

\gkanswer{答案（9）：$\dfrac{\sqrt5}{5}$}

\gkanswer{答案（10）：$\pi$}

\end{answer}
\begin{answer}
\gksolhead{2}{}

\textbf{题目}：

（本题满分 15 分）

本题共有 5 个小题，每一个小题都给出代号为 A、B、C、D 的四个结论，其中只有一个结论是正确的，把正确结论的代号写在题后的圆括号内，选对得 3 分，不选、选错或者选出的代号超过一个（不论是否都写在圆括号内），一律得 0 分。

（1） 设 $z$ 为复数，$z$ 是 $\overline z$ 的共轭复数，则 $z=\overline z$ 是 $z$ 为实数的

（A）充分但不必要的条件　　（B）必要但不充分的条件　　（C）充分而且必要的条件　　（D）既不充分也不必要的条件

（2） 下列各式中，正确的是

（A）$\arcsin\left(-\dfrac\pi3\right)=-\dfrac{\sqrt3}{2}$

（B）$\sin\left(\arcsin\dfrac\pi3\right)=\dfrac\pi3$

（C）$\arcsin\left(\sin\dfrac{5\pi}{4}\right)=-\dfrac\pi4$

（D）$\sin\left[\arccos\left(-\dfrac12\right)\right]=-\dfrac{\sqrt3}{2}$

\Needspace{5\baselineskip}
（3） 已知函数

\nopagebreak[4]
\[
f(x)=\lg(x^2-3x+2)
\]

\Needspace{5\baselineskip}
的定义域为 $F$，函数

\nopagebreak[4]
\[
g(x)=\lg(x-1)+\lg(x-2)
\]

的定义域为 $G$，那么

（A）$F\cap G=\varnothing$　　（B）$F=G$　　（C）$F\subset G$　　（D）$G\subset F$

（4） 若一个棱锥的底面是边数大于 3 的凸多边形，它的顶点到底面各边的距离都相等，则这个棱锥的底面多边形

（A）必为正多边形　　（B）必有内切圆　　（C）必有外接圆　　（D）必既有内切圆又有外接圆

（5） 若平移坐标轴，把坐标系 $xOy$ 的原点 $O$ 移到点 $O'$，$O'$ 在原坐标系中的坐标为 $(2,-1)$，则原坐标系中的曲线 $x=x^3$ 在新坐标系 $x'O'y'$ 中的方程是

（A）$y'+1=(x'-2)^3$　　（B）$y'+1=(x'+2)^3$

（C）$y'-1=(x'-2)^3$　　（D）$y'-1=(x'+2)^3$

\Needspace{4\baselineskip}
\textbf{解析}

\gksubpart{（1）}

复数满足 $z=\overline z$ 当且仅当其虚部为 0，也就是 $z$ 为实数。因此选（C）。

\gksubpart{（2）}

\Needspace{5\baselineskip}
$\arcsin$ 的定义域为 $[-1,1]$，所以（A）、（B）中的表达式均无定义。又

\nopagebreak[4]
\[
\arcsin\left(\sin\frac{5\pi}{4}\right)=\arcsin\left(-\frac{\sqrt2}{2}\right)=-\frac\pi4,
\]

故（C）正确；（D）中 $\sin(\arccos(-1/2))=\sqrt3/2$。因此选（C）。

\gksubpart{（3）}

\Needspace{5\baselineskip}
由 $x^2-3x+2=(x-1)(x-2)>0$ 得

\nopagebreak[4]
\[
F=(-\infty,1)\cup(2,+\infty).
\]

\Needspace{5\baselineskip}
而 $g(x)$ 有定义须同时满足 $x-1>0$、$x-2>0$，所以

\nopagebreak[4]
\[
G=(2,+\infty)\subset F.
\]

故选（D）。

\gksubpart{（4）}

设棱锥顶点在底面上的射影为 $O$。顶点到各底边的空间距离相等，意味着 $O$ 到底面各边所在直线的距离相等；对凸多边形的内侧取向，这些边都是同一圆的切线。因此底面多边形必有内切圆，选（B）。

\gksubpart{（5）}

\Needspace{5\baselineskip}
固定源字面为“原坐标系中的曲线 $x=x^3$”，而四个选项均为 $y=x^3$ 经原点平移到 $(2,-1)$ 后可能出现的形式。若按字面处理，曲线只涉及 $x$，平移后不会得到四个选项中的任何一个；若将左端改成 $y$，则由

\nopagebreak[4]
\[
x=x'+2,\ y=y'-1
\]

可得 $y'-1=(x'+2)^3$，即（D）。但后一结论依赖未经固定源支持的改字，故本题答案保留为未确定，并列为阻塞问题。

\gkanswer{答案（1）：（C）}

\gkanswer{答案（2）：（C）}

\gkanswer{答案（3）：（D）}

\gkanswer{答案（4）：（B）}

\gkanswer{答案（5）：固定源字面与选项冲突，暂不确定。}

\end{answer}
\begin{answer}
\gksolhead{3}{}

\textbf{题目}：

（本题满分 9 分）

本题共有 2 个小题，第（1）小题满分 5 分，第（2）小题满分 4 分。

\Needspace{5\baselineskip}
（1） 在直角坐标系内，方程

\nopagebreak[4]
\[
\left\{\begin{aligned}
x&=3\operatorname{ctg}\phi,\\
y&=2\operatorname{csc}\phi
\end{aligned}\right.
\]

（$\phi$ 是参数）表示什么曲线？画出它的图形。

\Needspace{5\baselineskip}
（2） 在极坐标系内，方程

\nopagebreak[4]
\[
\frac{2\sin\theta}{\rho}=1
\]

表示什么曲线？画出它的图形。
\nopagebreak[4]
\begin{center}% 1985 年上海卷理工农医类第 3 题：双曲线与圆
\begin{tikzpicture}[scale=0.72,>=stealth]
  \begin{scope}[shift={(-3.0,0)}]
    \draw[->] (-2.5,0)--(2.55,0) node[right] {$x$};
    \draw[->] (0,-2.7)--(0,2.7) node[above] {$y$};
    \draw[thick,domain=-2.2:2.2,samples=100]
      plot (\x,{2*sqrt(1+\x*\x/9)});
    \draw[thick,domain=-2.2:2.2,samples=100]
      plot (\x,{-2*sqrt(1+\x*\x/9)});
    \fill (0,2) circle (1.1pt);
    \fill (0,-2) circle (1.1pt);
    \node[below left] at (0,0) {$O$};
    \node[below] at (0,-3.0) {（1）};
  \end{scope}
  \begin{scope}[shift={(2.35,0)}]
    \draw[->] (-1.45,0)--(1.5,0) node[right] {$x$};
    \draw[->] (0,-0.45)--(0,2.55) node[above] {$y$};
    \draw[thick] (0,1) circle (1);
    \fill (0,1) circle (1.1pt);
    \node[below left] at (0,0) {$O$};
    \node[below] at (0,-0.75) {（2）};
  \end{scope}
\end{tikzpicture}
\end{center}
\nopagebreak[4]
\Needspace{4\baselineskip}
\textbf{解析}

\gksubpart{（1）}

\Needspace{5\baselineskip}
由 $x=3\operatorname{ctg}\phi$、$y=2\operatorname{csc}\phi$ 及恒等式 $\operatorname{csc}^2\phi-\operatorname{ctg}^2\phi=1$，得

\nopagebreak[4]
\[
\frac{y^2}{4}-\frac{x^2}{9}=1.
\]

这是中心在原点、焦点在 $y$ 轴上的双曲线，顶点为 $(0,\pm2)$。

\gksubpart{（2）}

\Needspace{5\baselineskip}
极坐标方程化为 $\rho=2\sin\theta$。两边乘以 $\rho$，利用 $x=\rho\cos\theta$、$y=\rho\sin\theta$，得

\nopagebreak[4]
\[
x^2+y^2=2y,
\]

\Needspace{5\baselineskip}
即

\nopagebreak[4]
\[
x^2+(y-1)^2=1.
\]

因此图形是圆心 $(0,1)$、半径为 $1$ 的圆。

\gkanswer{答案（1）：双曲线 $\dfrac{y^2}{4}-\dfrac{x^2}{9}=1$。}

\gkanswer{答案（2）：圆 $x^2+(y-1)^2=1$。}

\end{answer}
\begin{answer}
\gksolhead{4}{}

\textbf{题目}：

（本题满分 16 分）

本题共有 3 个小题，第（1）、（2）小题满分各 5 分，第（3）小题满分 6 分。

\Needspace{5\baselineskip}
（1） 求

\nopagebreak[4]
\[
\left(x^3-\frac1{x^2}\right)^{15}
\]

的展开式中的常数项。

（2） 从六个数字 $1,2,3,4,5,6$ 中任取四个不同的数字，有多少种取法？由这六个数字可以组成多少个没有重复数字的四位偶数？

（3） 在一列球中，第 1 个球的半径为 1，第 2 个球的直径是第 1 个球的半径，以后每一个球的直径都是前一个球的半径，这样无限继续下去，求所有这些球的表面积的和。

\Needspace{4\baselineskip}
\textbf{解析}

\gksubpart{（1）}

\Needspace{5\baselineskip}
通项中从第二项取 $k$ 次的项为

\nopagebreak[4]
\[
\binom{15}{k}x^{3(15-k)}\left(-\frac1{x^2}\right)^k
=\binom{15}{k}(-1)^k x^{45-5k}.
\]

\Needspace{5\baselineskip}
常数项要求 $45-5k=0$，故 $k=9$，常数项为

\nopagebreak[4]
\[
\binom{15}{9}(-1)^9=-5005.
\]

\gksubpart{（2）}

\Needspace{5\baselineskip}
从 6 个数字中任取 4 个，取法数为

\nopagebreak[4]
\[
\binom64=15.
\]

\Needspace{5\baselineskip}
四位偶数的个位有 $2\text{、}4\text{、}6$ 三种选择；确定个位后，百位、十位、千位依次有 $5\text{、}4\text{、}3$ 种选择。因此共有

\nopagebreak[4]
\[
3\cdot5\cdot4\cdot3=180
\]

个。

\gksubpart{（3）}

\Needspace{5\baselineskip}
第 $n$ 个球的半径为 $r_n=2^{-(n-1)}$，所以表面积总和为

\nopagebreak[4]
\[
\sum_{n=1}^{\infty}4\pi r_n^2
=4\pi\sum_{n=1}^{\infty}\left(\frac14\right)^{n-1}
=\frac{16\pi}{3}.
\]

\gkanswer{答案（1）：$-5005$}

\gkanswer{答案（2）：15 种取法，180 个四位偶数。}

\gkanswer{答案（3）：$\dfrac{16\pi}{3}$}

\end{answer}
\begin{answer}
\gksolhead{5}{}

\textbf{题目}：

（本题满分 6 分）

设 $\triangle ABC$ 的两个内角 $A\text{、}B$ 所对的边分别为 $a\text{、}b$。若复数 $z_1$ 的模为 $a$，幅角为 $B$；复数 $z_2$ 的模为 $b$，幅角为 $-A$。求证 $z_1+z_2$ 为实数。

\Needspace{4\baselineskip}
\textbf{解析}

\Needspace{5\baselineskip}
按复数的三角形式写出

\nopagebreak[4]
\[
z_1=a(\cos B+i\sin B),\ z_2=b(\cos A-i\sin A).
\]

\Needspace{5\baselineskip}
其和的虚部为 $a\sin B-b\sin A$。由正弦定理

\nopagebreak[4]
\[
\frac a{\sin A}=\frac b{\sin B},
\]

\Needspace{5\baselineskip}
可得 $a\sin B=b\sin A$，所以虚部为 0。因此

\nopagebreak[4]
\[
z_1+z_2\in\mathbb R.
\]

\gkanswer{答案：$z_1+z_2$ 为实数。}

\end{answer}
\begin{answer}
\gksolhead{6}{}

\textbf{题目}：

（本题满分 8 分）

\Needspace{5\baselineskip}
解三角方程

\nopagebreak[4]
\[
\cos^2x+\cos^23x=1
\]

\Needspace{4\baselineskip}
\textbf{解析}

\Needspace{5\baselineskip}
由二倍角公式，原方程等价于

\nopagebreak[4]
\[
\frac{1+\cos2x}{2}+\frac{1+\cos6x}{2}=1,
\]

\Needspace{5\baselineskip}
即

\nopagebreak[4]
\[
\cos2x+\cos6x=0
\Longleftrightarrow
2\cos4x\cos2x=0.
\]

\Needspace{5\baselineskip}
于是

\nopagebreak[4]
\[
\cos2x=0\Longleftrightarrow x=\frac\pi4+k\frac\pi2,
\]

\Needspace{5\baselineskip}
或

\nopagebreak[4]
\[
\cos4x=0\Longleftrightarrow x=\frac\pi8+k\frac\pi4,
k\in\mathbb Z.
\]

\gkanswer{答案：$x=\dfrac\pi4+k\dfrac\pi2$ 或 $x=\dfrac\pi8+k\dfrac\pi4$，$k\in\mathbb Z$。}

\end{answer}
\begin{answer}
\gksolhead{7}{}

\textbf{题目}：

（本题满分 6 分）

已知函数 $f(x)$ 对其定义域内的任意两个实数 $a\text{、}b$，当 $a<b$ 时，都有 $f(a)<f(b)$。试用反证法证明：方程 $f(x)=0$ 至多有一个实数根。

\Needspace{4\baselineskip}
\textbf{解析}

\Needspace{5\baselineskip}
反设方程 $f(x)=0$ 有两个不同的实数根 $x_1\text{、}x_2$。不妨设 $x_1<x_2$，由题设严格单调性应有

\nopagebreak[4]
\[
f(x_1)<f(x_2).
\]

但 $f(x_1)=f(x_2)=0$，矛盾。因此方程 $f(x)=0$ 至多有一个实数根。

\gkanswer{答案：方程 $f(x)=0$ 至多有一个实数根。}

\end{answer}
\begin{answer}
\gksolhead{8}{}

\textbf{题目}：

（本题满分 8 分）

\Needspace{5\baselineskip}
对于一切大于 1 的自然数 $n$，证明：

\nopagebreak[4]
\[
\left(1+\frac13\right)\left(1+\frac15\right)\cdots\left(1+\frac1{2n-1}\right)>\frac{\sqrt{2n+1}}2
\]

\Needspace{4\baselineskip}
\textbf{解析}

\Needspace{5\baselineskip}
记

\nopagebreak[4]
\[
P_n=\prod_{k=2}^{n}\left(1+\frac1{2k-1}\right).
\]

\Needspace{5\baselineskip}
当 $n=2$ 时，

\nopagebreak[4]
\[
P_2=\frac43>\frac{\sqrt5}{2}.
\]

\Needspace{5\baselineskip}
假设 $P_{n-1}>\frac{\sqrt{2n-1}}2$，则

\nopagebreak[4]
\[
P_n=P_{n-1}\frac{2n}{2n-1}
>\frac{\sqrt{2n-1}}2\cdot\frac{2n}{2n-1}.
\]

\Needspace{5\baselineskip}
又因为

\nopagebreak[4]
\[
\left(\frac{2n}{2n-1}\right)^2
>\frac{2n+1}{2n-1}
\]

\Needspace{5\baselineskip}
（等价于 $4n^2>(2n-1)(2n+1)$），所以

\nopagebreak[4]
\[
\frac{\sqrt{2n-1}}2\cdot\frac{2n}{2n-1}
>\frac{\sqrt{2n+1}}2.
\]

归纳可知题面不等式对一切 $n>1$ 成立。

\gkanswer{答案：按数学归纳法证明题面不等式成立。}

\end{answer}
\begin{answer}
\gksolhead{9}{}

\textbf{题目}：

（本题满分 10 分）

在三棱锥 $P-ABC$ 中，顶点 $P$ 到 $BC\text{、}AC\text{、}AB$ 的距离分别为 $h_1\text{、}h_2\text{、}h_3$，二面角 $P$- $BC$- $A=\alpha_1$，$P$- $AC$- $B=\alpha_2$，$P$- $AB$- $C=\alpha_3$。$\alpha_1\text{、}\alpha_2\text{、}\alpha_3$ 都是锐角，且依次成等差数列，$h_1\text{、}h_2\text{、}h_3$ 依次成等比数列。求证：$\alpha_1=\alpha_2=\alpha_3$。

\Needspace{4\baselineskip}
\textbf{解析}

\Needspace{5\baselineskip}
设点 $P$ 到底面 $ABC$ 的距离为 $H$。在分别垂直于 $BC\text{、}AC\text{、}AB$ 的截面中，二面角的平面角分别为 $\alpha_1\text{、}\alpha_2\text{、}\alpha_3$，故

\nopagebreak[4]
\[
H=h_1\sin\alpha_1=h_2\sin\alpha_2=h_3\sin\alpha_3.
\]

\Needspace{5\baselineskip}
由 $h_1\text{、}h_2\text{、}h_3$ 成等比数列，

\nopagebreak[4]
\[
h_2^2=h_1h_3.
\]

\Needspace{5\baselineskip}
代入 $h_i=H/\sin\alpha_i$，得

\nopagebreak[4]
\[
\sin^2\alpha_2=\sin\alpha_1\sin\alpha_3.
\]

\Needspace{5\baselineskip}
又因三角数列成等差数列，$\alpha_2=(\alpha_1+\alpha_3)/2$。利用恒等式

\nopagebreak[4]
\[
\sin^2\frac{\alpha_1+\alpha_3}{2}-\sin\alpha_1\sin\alpha_3
=\sin^2\frac{\alpha_1-\alpha_3}{2},
\]

\Needspace{5\baselineskip}
可知 $\alpha_1=\alpha_3$，从而

\nopagebreak[4]
\[
\alpha_1=\alpha_2=\alpha_3.
\]

\gkanswer{答案：$\alpha_1=\alpha_2=\alpha_3$。}

\end{answer}
\begin{answer}
\gksolhead{10}{}

\textbf{题目}：

（本题满分 12 分）

\Needspace{5\baselineskip}
已知：直线 $y=x+m$ 和曲线

\nopagebreak[4]
\[
x^2+2y^2+y^4-1=0
\]

交于 $A\text{、}B$ 两点，$P$ 是这条直线上的点，且 $|PA|\cdot|PB|=2$。求：当 $m$ 变化时，点 $P$ 的轨迹方程，并说明轨迹是什么图形。

\Needspace{4\baselineskip}
\textbf{解析}

\Needspace{5\baselineskip}
令 $P=(u,v)$。由于 $P$ 在直线 $y=x+m$ 上，故

\nopagebreak[4]
\[
m=v-u.
\]

\Needspace{5\baselineskip}
用交点的纵坐标 $z$ 作参数，直线上的点为 $(z-m,z)$。代入曲线方程得到交点纵坐标满足

\nopagebreak[4]
\[
h(z)=z^4+3z^2-2mz+m^2-1=0.
\]

\Needspace{5\baselineskip}
设 $A\text{、}B$ 的纵坐标为 $z_1\text{、}z_2$，令

\nopagebreak[4]
\[
s=z_1+z_2,\ r=z_1z_2.
\]

\Needspace{5\baselineskip}
另外两个根组成的二次因子记为 $z^2+sz+c$，于是

\nopagebreak[4]
\[
h(z)=(z^2-sz+r)(z^2+sz+c),
\]

\Needspace{5\baselineskip}
并且

\nopagebreak[4]
\[
c=3-r+s^2, s(c-r)=2m, rc=m^2-1.
\]

\Needspace{5\baselineskip}
由距离条件，直线方向向量为 $(1,1)$，所以

\nopagebreak[4]
\[
|PA|\cdot|PB|=2\,|(v-z_1)(v-z_2)|=2.
\]

\Needspace{5\baselineskip}
因此

\nopagebreak[4]
\[
v^2-sv+r=\varepsilon, \varepsilon\in\{1,-1\}.
\]

\Needspace{5\baselineskip}
将 $r=sv-v^2+\varepsilon$ 代入上面的系数关系并消去 $r,c$，得到

\nopagebreak[4]
\[
s^2=
\begin{cases}
u^2+v^4-3,&\varepsilon=1,\\
-(u^2+v^4+4v^2+3),&\varepsilon=-1.
\end{cases}
\]

\Needspace{5\baselineskip}
第二种情形右端严格为负，不可能。因此必有 $\varepsilon=1$。记

\nopagebreak[4]
\[
D=u^2+v^4-3=s^2.
\]

\Needspace{5\baselineskip}
在 $s(c-r)=2m$ 中代入 $r=sv-v^2+1$、$m=v-u$，并使用 $s^2=D$，可得

\nopagebreak[4]
\[
s(D+2v^2+1)=2(vD+v-u)=2\bigl[v(D+1)-u\bigr].
\]

\Needspace{5\baselineskip}
再平方消去 $s$，得到轨迹的代数方程

\nopagebreak[4]
\[
4\bigl[v(D+1)-u\bigr]^2=D\bigl(D+2v^2+1\bigr)^2,
D=u^2+v^4-3.
\]

\Needspace{5\baselineskip}
这是一个高次代数曲线，不是圆、椭圆或双曲线。由于题设要求直线与曲线有两个实交点，实际轨迹只取该方程上使

\nopagebreak[4]
\[
s=\frac{2\bigl[v(D+1)-u\bigr]}{D+2v^2+1}
\]

\Needspace{5\baselineskip}
对应的二次式 $z^2-sz+(sv-v^2+1)$ 有两个实根的部分，即

\nopagebreak[4]
\[
s^2-4sv+4v^2-4>0.
\]

若将相切的极限情形也计入，则将上式的 $>$ 改为 $\ge$。

\gkanswer{答案：见解析正文。}

\end{answer}